\scanpage{97}
\begin{center}- \mspage{156} -\end{center}
\margin{IX}
\[
(12.4.2)\qquad d(x_i)=C'_i-C''_i\quad,
\]
$C'_i$ et $C''_i$ étant les deux composantes irréductibles de $\widetilde C$ qui correspondent aux branches passant par le point double $x_i$. L'opérateur $d$ dépend du choix de l'ordre $(C'_i,C''_i)$ qu'on se donne entre ces deux branches, et l'isomorphisme (12.4.1) dépend évidemment de ce système de choix. Cependant, si pour $i\in I$ on désigne par
\[
(12.4.3)\qquad\varphi_{x_i}:M\longrightarrow\mathbb Z
\]
le composé $M\longrightarrow\mathbb Z^I\xrightarrow{\mathrm{pr}_i}\mathbb Z$, où le premier homomorphisme provient de (12.4.1), on constate aussitôt (en revenant à l'expression (12.3.4)) que $\varphi_{x_i}$ est bien déterminé au signe près ; de façon plus précise, il dépend seulement de l'ordre pris entre les deux branches $C'_i,C''_i$ passant par $i$, et change de signe quand on change cet ordre. Il s'ensuit notamment que la forme bilinéaire
\[
(12.4.4)\qquad\varphi_{x_i}\otimes\varphi_{x_i}\in\check M\otimes_{\mathbb Z}\check M\simeq\operatorname{Hom}(M\otimes M,\mathbb Z)
\]
est déterminée canoniquement en termes du point singulier $x_i$, sans ambiguïté de signe. On peut donc poser
\figure{\begin{tikzpicture}\draw[->] (0.7,0)--(0,0);\end{tikzpicture}}
\[
(12.4.5)\qquad u=\sum_i\varphi_{x_i}\otimes\varphi_{x_i}:M\otimes_{\mathbb Z}M\longrightarrow\mathbb Z\quad,
\]
et on trouve ainsi un accouplement bien déterminé i.e. indépendant de tout choix de signes. Il est clair qu'il est \uline{défini positif}, car c'est la restriction à $M\hookrightarrow\mathbb Z^I$ de la forme quadratique standart $\sum X_i^2$ sur $\mathbb Z^I$.

\margin{474}
\margin{474}
\margin{96}
